\section{Related Work}

\subsection{Slovene language resources and social-media corpora}

Written standard Slovene is well served by reference resources. Gigafida 2.0 \parencite{krek2020b-gf,gigafida-clarin} is the main reference corpus (1.2 billion tokens of newspaper, web, and book text, with automatic morphosyntactic annotation); ssj500k \parencite{krek2020a-ssj500k} provides gold-standard morphosyntactically annotated training data; siParl 3.0 \parencite{parliamentaryCorpus} covers parliamentary proceedings. The broader methodological context is surveyed in \textcite{studije2005korpusnem}.

The most directly relevant prior work is the JANES project (\textit{Jezikovna analiza nestandardne slovenščine}), which produced the first corpus of Slovene computer-mediated communication, drawing primarily on Twitter but also covering news comments and forum posts \parencite{erjavec2019janes}. JANES demonstrated both the feasibility of building Slovene social-media corpora and the specific challenges they pose: non-standard spelling, code-switching, short texts, and unreliable platform language tags. The present corpus continues in this tradition, shifting the target platform from Twitter to Bluesky and dealing with a community that has formed largely as a direct consequence of Twitter's decline.

\subsection{Bluesky and ATProto}

The migration of users from Twitter/X to alternative platforms accelerated sharply after Elon Musk's acquisition of Twitter in late 2022 and intensified further around the 2024 US presidential election. Bluesky became the main destination for a segment of this migration --- partly because its interface is close to what Twitter users were familiar with, making it considerably more accessible than federated alternatives such as Mastodon, which require users to choose a server instance and navigate inter-instance social friction. The corpus collected for this study directly captures the consequences of this shift, visible as a sharp late-2024 increase in Slovene post volume and a concentration of authors' first included posts in late 2024 (Section~\ref{sec:empirical}).

Technically, Bluesky is built on the AT Protocol (ATProto), a federated specification that separates identity from hosting: user identities are globally portable DIDs (decentralised identifiers) resolved through a public directory, and posts are stored on independently operated Personal Data Servers \parencite{atproto-spec}. From a corpus-collection standpoint, ATProto has one decisive advantage over proprietary platforms: all public content is accessible via a network-wide event firehose and per-account repository endpoints, without API quotas or paywalls. This makes collection reproducible at low cost, and the DID resolution mechanism allows a collector to follow an account regardless of which PDS hosts it --- a property exploited by the multi-PDS collection approach described in this paper.

Existing empirical work on Bluesky has so far operated at large scale and focused primarily on English-language content and network-level phenomena: \textcite{failla2024bluesky} assembled a dataset of 235 million posts from over four million users for platform-wide behavioural analysis, while \textcite{quelle2025bluesky} used Bluesky data to characterise network topology and political polarization. Language-specific corpus construction for a smaller community on ATProto has not previously been reported in the literature.

\subsection{Language identification in corpus construction}

Building a language-specific corpus from a multilingual stream typically combines two signals: platform-supplied language metadata and automatic language identification. Neither is reliable on its own. Platform metadata is user- or app-generated: in the present corpus, 41.4\% of posts carry both \texttt{sl} and \texttt{en} tags, almost entirely because Bluesky's default interface language is English, not because the posts are bilingual. Automatic language identification tools---langid.py \parencite{lui2012langid} and langdetect \parencite{nakatani2010langdetect} are used here---perform better at post level, but both struggle with short texts and with closely related languages; Slovene, Croatian, and Bosnian are routinely confused by standard detectors. Using consensus across multiple detectors and requiring agreement with author-supplied tags improves precision, which is the rationale for the tiered inclusion strategy described in Section~\ref{sec:corpus-construction}.
